\section*{Partie 1 — Modélisation orientée objet}
\subsection*{Objectif}
Dans cette première partie, vous allez proposer une modélisation orientée objet des composants du système de palettisation. L’objectif est d’identifier les objets, leurs missions, leurs états et leurs interactions.

\subsection*{Question 1 — Identification des objets (4 points)}
Listez les objets fonctionnels présents dans le système de palettisation automatique. Pour chacun, indiquez :
\begin{itemize}
    \item Un court descriptif de son rôle dans le système
    \item Deux missions possibles qu’il peut recevoir
    \item Trois états ou propriétés qu’il peut exposer
\end{itemize}

\subsection*{Question 2 — Structure d'état (2 points)}
Choisissez un des objets listés en Question 1 (autre que le superviseur) et proposez une structure `STRUCT` décrivant son état. Elle devra comporter au minimum trois champs pertinents (booléens, entiers, énumérations, etc.).

\subsection*{Question 3 — Interactions entre objets (2 points)}
Expliquez comment les objets suivants interagissent entre eux pendant un cycle complet de constitution d'une couche de boîtes dans la zone tampon :
\begin{itemize}
    \item Le superviseur
    \item La station d’alimentation
    \item Le robot palettiseur
\end{itemize}
Vous pouvez accompagner votre réponse d’un schéma simple (texte ou croquis).

\subsection*{Barème indicatif}
\begin{itemize}
    \item Q1 : /4 — identification complète, missions pertinentes, états cohérents
    \item Q2 : /2 — structure adaptée, noms explicites, typage correct
    \item Q3 : /2 — logique claire, interactions bien expliquées
\end{itemize}
